\documentclass[12pt]{article}

% Preamble
\usepackage{amsmath, amssymb, graphicx, color}
\usepackage{hyperref}
\usepackage{geometry}

% Set page margins, etc.
\geometry{a4paper, margin=1in}

% Document metadata
\title{Global Regularity Framework for the Navier-Stokes Equations}
\author{Your Name}
\date{\today}

% Begin the document
\begin{document}

% Title page
\maketitle

% Abstract
\begin{abstract}
    This document consolidates various research findings related to the Millennium Prize Problem, focusing on the global regularity of the Navier-Stokes equations. Through the integration of numerous theoretical frameworks, a comprehensive approach to understanding the behavior of solutions is proposed. Each section presents critical aspects of the theory, aiming to contribute to a clearer understanding of the mathematical properties underlying the 3D Navier-Stokes equations.
\end{abstract}

% Table of Contents (optional, can be uncommented if needed)
%\tableofcontents

% Section: Introduction
\section{Introduction}
This document compiles and refines the research findings in a cohesive narrative. The global regularity problem for the 3D Navier-Stokes equations is central to the Millennium Prize Problem, and this document brings together various frameworks for understanding the solution properties. The following sections are drawn from individual works, each contributing a unique perspective to the overall problem.

% Section: Compilation of Documents
\section{Compilation of Documents}

% Add each file into the main document with the \input command.
% Ensure all .tex files are located in the same directory as main.tex or adjust paths accordingly.

\input{project-documents/Navier-Stokes/A_Unified_Framework_for_the_Global_Regularity_of_3D_Navier-Stokes_EquationsMoreover,_the_framework_offers_a_transferable_methodology_for_stabilizing_other_complex_PDE_system.tex}  

% Conclusion: Final Remarks
\section{Conclusion}
This compilation provides a thorough examination of the global regularity problem for the 3D Navier-Stokes equations, aggregating several significant frameworks and approaches. The presented work serves as a comprehensive reference and theoretical guide, contributing to the ongoing mathematical and computational efforts aimed at resolving the Millennium Prize Problem.

% End of the document
\end{document}